\subsubsection{Systém priateľstiev}

Systém priateľstiev umožňuje používateľom nadviazať spojenie a komunikovať
priamo cez správy. Implementácia pokrýva posielanie žiadostí, ich
prijímanie a odmietanie, odstraňovanie priateľov a priamu komunikáciu.

\paragraph{Dátový model:}
Každé priateľstvo je reprezentované triedou \texttt{Friendship} a uložené
do Firestore kolekcie \texttt{friendships}:

\begin{lstlisting}[language=Kotlin, caption=Dátová trieda Friendship]
data class Friendship(
    val id         : String = "",
    val requesterId: String = "",        // kto poslal žiadosť
    val receiverId : String = "",        // kto prijíma žiadosť
    val status     : String = "pending", // "pending" | "accepted"
    val timestamp  : Long   = 0L
)
\end{lstlisting}

\paragraph{Stavy priateľstva:}
Pre každú dvojicu používateľov môže existovať práve jeden vzťah.
Metóda \texttt{getFriendshipStatus()} vráti jeden zo štyroch stavov
uvedených v tabuľke~\ref{tab:friendship-states}:

\begin{lstlisting}[language=Kotlin, caption=Určenie stavu priateľstva]
fun getFriendshipStatus(currentUserId: String, otherUserId: String): String {
    val f = _friendships.firstOrNull {
        (it.requesterId == currentUserId && it.receiverId == otherUserId) ||
        (it.requesterId == otherUserId   && it.receiverId == currentUserId)
    } ?: return "none"

    return when {
        f.status == "accepted"         -> "friends"
        f.requesterId == currentUserId -> "pending_sent"
        else                           -> "pending_received"
    }
}
\end{lstlisting}

\begin{table}[h]
    \centering
    \caption{Stavy priateľstva}
    \label{tab:friendship-states}
    \begin{tabular}{ll}
        \toprule
        \textbf{Stav} & \textbf{Význam} \\
        \midrule
        \texttt{none}             & Žiadny vzťah medzi používateľmi. \\
        \texttt{pending\_sent}     & Aktuálny používateľ poslal žiadosť. \\
        \texttt{pending\_received} & Aktuálny používateľ dostal žiadosť. \\
        \texttt{friends}           & Priateľstvo je akceptované. \\
        \bottomrule
    \end{tabular}
\end{table}

Stav sa využíva v \texttt{PublicProfileScreen} na zobrazenie správneho
tlačidla — „Pridať priateľa", „Žiadosť odoslaná", „Prijať\,/\,Odmietnuť"
alebo „Odstrániť priateľa".

\paragraph{Načítanie priateľstiev:}
Pri prihlásení sa načítajú všetky dokumenty, kde je používateľ buď
odosielateľom alebo prijímateľom. Odmietnuté žiadosti sa nenačítavajú —
zabraňuje sa tým opakovanému posielaniu žiadostí od toho istého
používateľa:

\begin{lstlisting}[language=Kotlin, caption=Načítanie priateľstiev z Firestore]
fun loadFriendships(uid: String) {
    val sent     = db.collection("friendships")
        .whereEqualTo("requesterId", uid).get().await()
    val received = db.collection("friendships")
        .whereEqualTo("receiverId",  uid).get().await()
    (sent.documents + received.documents).forEach { doc ->
        val f = Friendship.fromFirestore(doc.id, doc.data!!)
        if (f.status != "declined") _friendships.add(f)
    }
}
\end{lstlisting}

\paragraph{Odoslanie žiadosti:}
Pred vytvorením novej žiadosti sa skontroluje, či medzi dvojicou
používateľov už vzťah existuje:

\begin{lstlisting}[language=Kotlin, caption=Odoslanie žiadosti o priateľstvo]
fun sendFriendRequest(fromId: String, toId: String) {
    val alreadyExists = _friendships.any {
        (it.requesterId == fromId && it.receiverId == toId) ||
        (it.requesterId == toId   && it.receiverId == fromId)
    }
    if (alreadyExists) {
        friendshipError = "Žiadosť o priateľstvo už existuje."
        return
    }
    val friendship = Friendship(
        id          = docRef.id,
        requesterId = fromId,
        receiverId  = toId,
        status      = "pending",
        timestamp   = System.currentTimeMillis()
    )
    docRef.set(friendship.toMap()).await()
    _friendships.add(friendship)
}
\end{lstlisting}

\paragraph{Odpoveď na žiadosť:}
Pri prijatí sa stav aktualizuje na \texttt{accepted}; pri odmietnutí
sa dokument z Firestore aj z lokálnej pamäte vymaže:

\begin{lstlisting}[language=Kotlin, caption=Prijatie alebo odmietnutie žiadosti]
fun respondToFriendRequest(friendshipId: String, accept: Boolean) {
    if (accept) {
        db.collection("friendships").document(friendshipId)
            .update("status", "accepted").await()
        _friendships[idx] = _friendships[idx].copy(status = "accepted")
    } else {
        db.collection("friendships").document(friendshipId)
            .delete().await()
        _friendships.removeAll { it.id == friendshipId }
    }
}
\end{lstlisting}

\paragraph{Odstránenie priateľa:}
Metóda \texttt{removeFriend()} vyhľadá akceptované priateľstvo medzi
dvojicou a vymaže ho z Firestore aj z lokálnej pamäte:

\begin{lstlisting}[language=Kotlin, caption=Odstránenie priateľa]
fun removeFriend(currentUserId: String, otherUserId: String) {
    val f = _friendships.firstOrNull {
        it.status == "accepted" &&
        ((it.requesterId == currentUserId && it.receiverId == otherUserId) ||
         (it.requesterId == otherUserId   && it.receiverId == currentUserId))
    } ?: return
    _friendships.remove(f)
    db.collection("friendships").document(f.id).delete().await()
}
\end{lstlisting}

\paragraph{Priama komunikácia:}
Priatelia si môžu písať správy cez \texttt{FriendChatScreen}. ID
konverzácie sa generuje deterministicky zo zoradených UID oboch
používateľov — zaručuje sa tak, že obaja vždy pristupujú k tomu istému
dokumentu bez ohľadu na to, kto rozhovor začal:

\begin{lstlisting}[language=Kotlin, caption=Deterministické generovanie ID konverzácie]
private fun directConversationId(uid1: String, uid2: String) =
    if (uid1 < uid2) "${uid1}_${uid2}" else "${uid2}_${uid1}"
\end{lstlisting}

Správy sa ukladajú do Firestore kolekcie
\texttt{directMessages/\{conversationId\}/messages} a synchronizujú sa
v reálnom čase cez \texttt{snapshotListener} — rovnaký mechanizmus
ako chat v udalostiach.

\paragraph{Používateľské rozhranie:}
Obrazovka \texttt{FriendsScreen} je dostupná z bočného menu pod položkou
„Priatelia" a je rozdelená do dvoch sekcií:
\begin{itemize}
    \item \textbf{Žiadosti o priateľstvo} — zobrazujú sa len ak existujú,
          s tlačidlami „Prijať" a „Odmietnuť".
    \item \textbf{Priatelia} — zoznam akceptovaných priateľov s možnosťou
          otvoriť priamu konverzáciu alebo priateľa odstrániť.
\end{itemize}